% p397_1; source PDF page 397; manually verified editable table.
\begin{tikzpicture}[x=1cm,y=-1cm,inner sep=0pt,outer sep=0pt]
\path[use as bounding box] (0,0) rectangle (16.000,3.520);
\draw[line width=.4pt] (0.000,0.000) rectangle (16.000,0.805);
\node[align=center,text width=15.740cm,font={\linespread{1}\fontsize{10}{13.50}\selectfont\bfseries}] at (8.000,0.403) {贪污罪=A+B};
\draw[line width=.4pt] (0.000,0.805) rectangle (5.300,3.520);
\node[align=center,text width=5.040cm,font={\linespread{1}\fontsize{10}{13.50}\selectfont}] at (2.650,2.163) {A：实行行为\\（侵占罪、盗窃罪、诈骗罪）};
\draw[line width=.4pt] (5.300,0.805) rectangle (16.000,1.710);
\node[align=left,text width=10.440cm,font={\linespread{1}\fontsize{10}{13.50}\selectfont}] at (10.650,1.258) {B1：行为主体（国家工作人员）};
\draw[line width=.4pt] (5.300,1.710) rectangle (16.000,2.615);
\node[align=left,text width=10.440cm,font={\linespread{1}\fontsize{10}{13.50}\selectfont}] at (10.650,2.163) {B2：行为对象（公共财物，入罪数额是 3 万元）};
\draw[line width=.4pt] (5.300,2.615) rectangle (16.000,3.520);
\node[align=left,text width=10.440cm,font={\linespread{1}\fontsize{10}{13.50}\selectfont}] at (10.650,3.068) {B3：行为方式（利用职务便利）};
\end{tikzpicture}
